% =========================================================
% Square Bracketing and Nonlinear Approximation
% =========================================================

\section{Square Bracketing and Nonlinear Approximation}
\label{sec:square-bracketing-nonlinear-approximation}

Linear approximation controls distances on the number line. Nonlinear
approximation asks for rational inputs whose images under nonlinear
operations, such as squaring, bracket or approximate a target value.

\begin{topicbox}{Square Approximation Toolkit}
\begin{exposition}
The basic square-approximation strategy is to trap a target positive rational
$r$ between two rational squares, $a^2 < r < b^2$, and then refine the upper
bracket using a fine dyadic grid. The least grid point whose square crosses $r$
gives an approximation from above with controlled overshoot.
\end{exposition}

\begin{tcolorbox}[colback=thmbox, colframe=thmborder, arc=2pt,
  left=6pt, right=6pt, top=4pt, bottom=4pt,
  title={\small\textbf{Theorem (Rational Square Bracketing)}},
  fonttitle=\small\bfseries]
\begin{theorem}[Rational Square Bracketing]
\label{thm:rational-square-bracketing}
For every $r\in\mathbb{Q}$ with $r>0$, there exist $a,b\in\mathbb{Q}$ with
$a>0$ and $b>0$ such that
\[
a^2<r<b^2.
\]
\hyperref[prf:rational-square-bracketing]{Go to proof.}
\end{theorem}
\end{tcolorbox}

\begin{remark*}[Standard quantified statement]
\[
\forall r\in\mathbb{Q},\;r>0\;\;
\exists a,b\in\mathbb{Q},\;a>0,\;b>0\;\;
a^2 < r < b^2.
\]
\end{remark*}

\begin{remark*}[Negated quantified statement]
\[
\exists r\in\mathbb{Q},\;r>0\;\;
\forall a,b\in\mathbb{Q},\;a>0,\;b>0\;\;
r\le a^2 \lor b^2\le r.
\]
\end{remark*}

\begin{remark*}[Interpretation]
Every positive rational lies strictly between the squares of two positive
rationals. The witnesses are $a = r/(r+1)$ and $b = r+1$: one checks
$a^2 < r$ and $r < b^2$ directly from the ordered field arithmetic of
$\mathbb{Q}$.
\end{remark*}

\begin{remark*}[Dependencies]
\begin{itemize}
  \item \hyperref[thm:q-is-ordered-field]{$\mathbb{Q}$ Is an Ordered Field}
  \item \hyperref[thm:archimedean-q]{Archimedean Property of $\mathbb{Q}$}
\end{itemize}
\end{remark*}
\end{topicbox}

\begin{topicbox}{Grid Crossing and Approximation from Above}
\begin{exposition}
Given a rational target $r > 0$ and a grid of step size $1/q$, the well-ordering
principle guarantees a least integer $p$ such that $(p/q)^2 > r$. The
predecessor $(p-1)/q$ has square at most $r$, and the overshoot
$(p/q)^2 - r$ is controlled by $2p/q^2$, which can be made smaller than
any prescribed $1/n$ by taking $q$ large enough.
\end{exposition}

\begin{lemma}[First Grid Point Crossing]
\label{lem:first-grid-point-crossing}
Let $r\in\mathbb{Q}$ with $r>0$ and let $q\in\mathbb{N}$ with $q\ge 1$.
If $K\in\mathbb{N}$ satisfies $K^2>r$, then the set
\[
A=\{m\in\mathbb{N}:m^2>rq^2\}
\]
is nonempty and therefore has a least element.
\hyperref[prf:first-grid-point-crossing]{Go to proof.}
\end{lemma}

\begin{remark*}[Standard quantified statement]
\[
\forall r\in\mathbb{Q},\;r>0\;\;
\forall q\in\mathbb{N},\;q\ge1\;\;
\bigl(\exists K\in\mathbb{N}\;(K^2>r)\bigr)
\Longrightarrow
\exists p_0\in\mathbb{N}\;\bigl(p_0\in A \land \forall m\in A\;(p_0\le m)\bigr).
\]
\end{remark*}

\begin{remark*}[Interpretation]
The set $A$ is nonempty because $Kq \in A$. By the Well-Ordering Principle,
every nonempty subset of $\mathbb{N}$ has a least element. That least element
$p$ is the first grid point $p/q$ whose square exceeds $r$.
\end{remark*}

\begin{remark*}[Dependencies]
\begin{itemize}
  \item \hyperref[thm:well-ordering-principle]{Well-Ordering Principle}
  \item \hyperref[thm:rational-square-bracketing]{Rational Square Bracketing}
\end{itemize}
\end{remark*}

\begin{tcolorbox}[colback=thmbox, colframe=thmborder, arc=2pt,
  left=6pt, right=6pt, top=4pt, bottom=4pt,
  title={\small\textbf{Theorem (Rational Square Approximation from Above)}},
  fonttitle=\small\bfseries]
\begin{theorem}[Rational Square Approximation from Above]
\label{thm:rational-square-approximation-from-above}
Let $r\in\mathbb{Q}$ with $r>0$, and let $n\in\mathbb{N}$ with $n\ge 1$.
Then there exists $s\in\mathbb{Q}$ with $s>0$ such that
\[
r<s^2<r+\frac1n.
\]
\hyperref[prf:rational-square-approximation-from-above]{Go to proof.}
\end{theorem}
\end{tcolorbox}

\begin{remark*}[Standard quantified statement]
\[
\forall r\in\mathbb{Q},\;r>0\;\;
\forall n\in\mathbb{N},\;n\ge1\;\;
\exists s\in\mathbb{Q},\;s>0\;\;
r < s^2 < r+\tfrac{1}{n}.
\]
\end{remark*}

\begin{remark*}[Interpretation]
Choose grid denominator $q$ large enough that $2Kq^{-1} < 1/n$ where $K$ is a
rational upper bound for $\sqrt{r}$. Let $p$ be the least natural number with
$p^2 > rq^2$ and set $s = p/q$. Then $s^2 > r$ by construction and
$s^2 - r = (p^2 - rq^2)/q^2 \le (2p-1)/q^2 < 2K/q < 1/n$.
\end{remark*}

\begin{remark*}[Dependencies]
\begin{itemize}
  \item \hyperref[lem:first-grid-point-crossing]{First Grid Point Crossing}
  \item \hyperref[thm:archimedean-q]{Archimedean Property of $\mathbb{Q}$}
  \item \hyperref[thm:rational-square-bracketing]{Rational Square Bracketing}
\end{itemize}
\end{remark*}

\begin{tcolorbox}[colback=thmbox, colframe=thmborder, arc=2pt,
  left=6pt, right=6pt, top=4pt, bottom=4pt,
  title={\small\textbf{Theorem (Rational Square Approximation from Below)}},
  fonttitle=\small\bfseries]
\begin{theorem}[Rational Square Approximation from Below]
\label{thm:rational-square-approximation-from-below}
Let $r\in\mathbb{Q}$ with $r>0$, and let $n\in\mathbb{N}$ with $n\ge 1$.
If $r-1/n>0$, then there exists $s\in\mathbb{Q}$ with $s>0$ such that
\[
r-\frac1n<s^2<r.
\]
\hyperref[prf:rational-square-approximation-from-below]{Go to proof.}
\end{theorem}
\end{tcolorbox}

\begin{remark*}[Standard quantified statement]
\[
\forall r\in\mathbb{Q},\;r>0\;\;
\forall n\in\mathbb{N},\;n\ge1,\;r-\tfrac{1}{n}>0\;\;
\exists s\in\mathbb{Q},\;s>0\;\;
r - \tfrac{1}{n} < s^2 < r.
\]
\end{remark*}

\begin{remark*}[Interpretation]
The rational $r - 1/n$ is positive by hypothesis. By the Approximation from
Above theorem applied to $r - 1/n$, there exists $s$ with
$(r-1/n) < s^2 < (r-1/n) + 1/n = r$.
\end{remark*}

\begin{remark*}[Dependencies]
\begin{itemize}
  \item \hyperref[thm:rational-square-approximation-from-above]{Rational Square Approximation from Above}
\end{itemize}
\end{remark*}
\end{topicbox}
